% ============================================================
%  CONCLUSION GÉNÉRALE
% ============================================================
\chapter*{Conclusion générale et perspectives}
\addcontentsline{toc}{chapter}{Conclusion générale et perspectives}
\markboth{Conclusion générale}{Conclusion générale}

\hspace{1cm}Ce projet a porté sur l'optimisation de la gestion des stocks au
sein de la \gls{sc} du groupe \gls{ocp}. Il s'agissait de dimensionner les
règles de réapprovisionnement de manière à répondre à la demande de
l'exploitation tout en maîtrisant le capital immobilisé.

\vspace{0.5cm}
\textbf{Bilan des réalisations :}

\begin{itemize}
  \item Une chaîne de traitement reproductible a été mise en place, de l'export
    brut jusqu'aux paramètres de gestion, à partir d'outils exclusivement
    libres.
  \item Le portefeuille de références a été segmenté par la matrice
    \gls{abc}/\gls{xyz}, permettant d'associer à chaque segment un mode de
    gestion adapté plutôt qu'une règle unique.
  \item Les stocks de sécurité et points de commande ont été recalculés en
    tenant compte simultanément de la variabilité de la demande et de celle des
    délais fournisseurs.
  \item [Rappeler ici les gains chiffrés obtenus, en taux de service et en
    valeur de stock.]
  \item Un tableau de bord a été livré aux équipes, leur permettant de simuler
    l'effet d'un taux de service cible avant décision.
\end{itemize}

\vspace{0.5cm}
\textbf{Limites du travail :}

Les résultats reposent sur une simulation à partir de l'historique disponible.
Ils dépendent d'hypothèses de coûts qui devront être confirmées, et l'hypothèse
de normalité de la demande reste discutable pour les références à consommation
intermittente.

\vspace{0.5cm}
\textbf{Perspectives :}

\begin{enumerate}
  \item Traiter les références à demande intermittente par des méthodes
    dédiées (Croston, lois de Poisson) plutôt que par l'approximation normale.
  \item Automatiser l'actualisation des paramètres par une connexion directe à
    l'\gls{erp}, en remplacement de l'export manuel.
  \item Étendre la démarche à l'amont de la chaîne : sélection et évaluation des
    fournisseurs sur la fiabilité de leurs délais, levier plus efficace que
    l'augmentation du stock.
  \item Mettre en place un suivi périodique des indicateurs afin de mesurer
    l'écart entre les gains simulés et les gains réellement constatés.
  \item Élargir le périmètre aux autres familles d'articles et aux autres sites
    du groupe.
\end{enumerate}
